
\subsubsection{Benchmark Taxonomies}

Papers with Code and similar platforms organize benchmarks by task type (classification, generation, translation) and domain (vision, language, multimodal). These taxonomies facilitate browsing but require manual review to determine hypothesis-specific suitability. The present work automates this matching via design feature extraction and historical validation.

\subsubsection{Citation-Based Analysis}

Semantic Scholar and Google Scholar enable benchmark popularity analysis through citation counts and co-citation networks. However, popularity does not equal coverage: highly cited benchmarks may be unsuitable for specific hypothesis types due to metric incompatibility. This work extracts design features to predict benchmark-hypothesis compatibility independent of popularity.

\subsubsection{Benchmark Recommendation Systems}

Collaborative filtering and user-behavior similarity approaches recommend benchmarks based on historical usage patterns. These methods require prior usage data for each hypothesis type, limiting applicability to novel categories. Design-feature-based approaches enable prediction without prior usage by leveraging inherent benchmark characteristics.

\subsubsection{Technical Foundations}

SciBERT provides scientific text understanding for citation context classification. SentenceBERT embeddings capture semantic similarity enabling clustering without pre-defined categories. K-means clustering provides unsupervised discovery of coverage families, complemented by silhouette scoring for cluster quality validation.
